\section{Discussion}
\label{sec:discussion}
 
\subsection{Causal interpretation and economic significance}
\label{sec:disc-causal}
 
This paper interprets its central estimate causally: subject to the identification assumptions defended in Section~\ref{sec:identification-assumptions} -- correct treatment allocation, no anticipation, and no spillovers -- together with the parallel-trends evidence in Section~\ref{sec:parallel-trends}, the Scotland$\times$Post coefficient is read as SCP's intention-to-treat effect on child material deprivation, not merely an association. The specification ladder assembled in Section~\ref{sec:res-dml} supports this reading: sign and rough order of magnitude survive every change to functional form and covariate breadth. If the OLS estimate were an artefact of a misspecified linear relationship, replacing it with lasso or random forest -- which share no functional-form assumption with OLS or each other -- would be expected to produce noisy or sign-inconsistent results, not a still-negative, still-clustered estimate. The OLS-to-DML gap (Section~\ref{sec:res-dml}) is nonetheless informative in its own right: OLS's covariate adjustment overstates the effect relative to every DML cell, consistent with some genuine nonlinearity in how the six CASE controls relate to deprivation that a linear model cannot capture. The DML models are also cross-fitted on a modest post-treatment Scottish sample (around 800 observations); refitting the preferred wide-covariate random-forest specification 20 times under different cross-fitting seeds finds the coefficient consistently negative but ranging from $-0.021$ to $-0.053$, clearing conventional 5\% significance in only a quarter of refits (Appendix Table~\ref{tab:dmlstability}).
 
The Adjusted OLS estimate implies material deprivation fell by 7.6 percentage points across Scotland's child population as a result of SCP; applied to Scotland's roughly 900,000 children under 16 (16.2 per cent of Scotland's 5.55 million population; \citep{nrs2025}), this is approximately 68,000 children. The doubly robust random-forest estimate under the wide covariate set ($-4.5$ percentage points) implies a smaller figure of around 40,500. This estimate is preferred here over the lean-covariate random-forest cell ($-4.0$ points) primarily because it tests sensitivity to the fuller covariate set rather than just the six CASE controls (Section~\ref{sec:res-dml}); its edge in reported significance over the lean cell should be read cautiously, since the stability check above finds that specific significance level holds in only a minority of refits. Both are back-of-envelope translations of a percentage-point coefficient into a headcount, not administrative estimates. Both also describe the intention-to-treat effect on the whole Scottish child population (Section~\ref{sec:identification-assumptions}), including non-recipient households whose deprivation should not move mechanically, so they likely understate the effect actually experienced by SCP recipients.
 
The Scottish Government's own modelled estimate of around 40,000 children kept out of relative poverty by SCP in 2025--26 offers one reference point \citep{scottishgovernment2025} (Section~\ref{sec:introduction}), but it describes income poverty, a related but distinct construct (Section~\ref{sec:lit-measurement}), so the comparison is only suggestive. A closer benchmark comes from \citet{andersen2025}'s own Table 3 estimates for the same outcomes ($-0.085$ for the MDCH flag, $-0.088$ for food insecurity). This paper's OLS Adjusted estimate ($-0.076$) sits just below that benchmark, and the DML estimates sit further below, though the comparison is not an independent check, since both papers draw on the same FRS/HBAI data.
 
\subsection{Reconciling design tensions}
\label{sec:disc-reconciliation}
 
The quarterly parallel-trends check (Section~\ref{sec:parallel-trends}) flags two individual items, Bed/bedroom and Outdoor play area, that pass annual-level testing but not quarterly. Bed/bedroom is asked only under the pre-FYE2024 question wording and is estimated on a sample roughly a sixth the size of most other items (Table~\ref{tab:stage1_items}); splitting an already-thin sample into 13-week windows leaves individual pre-treatment quarters imprecisely estimated, so a single noisy quarter crossing significance is unsurprising and does not indicate a real differential trend. Outdoor play area is less easily explained this way: its sample is not thin ($N=46{,}924$) and its annual-level coefficient is small and far from significant ($-0.013$, raw $p=0.499$), so a quarterly-only flag most plausibly reflects a single noisy pre-treatment quarter rather than a genuine trend, though this cannot be confirmed and remains an unresolved limitation. One of the two composite outcomes examined as a robustness check (Appendix~\ref{sec:appendix-composites}) is also flagged at quarterly level, alongside an unexplained placebo anomaly on Northern Ireland; since neither involves the headline official MDCH flag or the item-level results, they are discussed in full in the appendix rather than here.
 
Section~\ref{sec:identification-assumptions} carried forward one spillover threat that could not be tested directly: sustained UK-wide attention to SCP shifting England's own counterfactual trajectory through UK Government policy. The clearest such channel identified in this paper is the UK Government's removal of the two-child Universal Credit limit, which the Joseph Rowntree Foundation credits partly to Scotland's example (Section~\ref{sec:introduction}); the change takes effect from April 2026, after this paper's sample window ends (FYE2024), so it cannot itself have contaminated the English control trend behind the results reported here. This does not rule out lower-level media or political spillover earlier in the sample period, which remains untestable within this design and is carried forward as an honest limitation rather than a resolved one.
 
The item-level null result (Section~\ref{sec:res-item-dml}) counsels caution against reading the aggregate effect as evidence of a specific mechanism: no item carries a robust, correctable signal under any doubly robust re-estimation, and even the two items significant under simple OLS are not robust to whether the ten items are modelled independently or jointly (Section~\ref{sec:res-stacked}). Celebrations holds under both specifications, but Warm coat survives only the independent-regression one; unlike the stacked model, that specification does not account for the same child answering all ten items, understating its true standard error in exactly the way \citet{moulton1990} predicts, so the stacked model's null result for Warm coat is the more credible of the two rather than a merely competing reading. Combined with the aggregate composite's own consistency across estimators, the most defensible reading is that SCP moved material circumstances broadly, in ways too diffuse across the ten measured hardships for any single item to carry a robust signal on its own, rather than that SCP left specific hardships untouched. Separately, since SCP is not Scotland's only Scotland-only child transfer, some of the attributed effect may reflect Best Start Grant and Best Start Foods rather than SCP in isolation (Section~\ref{sec:identification-assumptions}); the age-restricted robustness check (Section~\ref{sec:res-robustness}) is reassuring but does not fully separate the two, since Best Start payments continue, at a lower rate, past age six.
 
This diffuseness has a plausible economic reading. SCP is unrestricted cash, not a transfer earmarked for a specific hardship, so standard consumer theory predicts it relaxes a household's whole budget constraint rather than compelling spending on one good; which constraint binds tightest, and which item improves as a result, should vary across recipients. \citet{andersen2025}'s own qualitative interviews with SCP recipients support this reading. Some parents described ringfencing the payment for child-specific items such as clothing, shoes and activities, while others pooled it into the household's general food and utility budget, a split driven by family size and the severity of existing hardship rather than a shared priority. \citet{nicoriciuelliot2025}'s latent class analysis of the same HBAI deprivation items suggests the same mechanism from the demand side. Households cluster into typologies defined not by how deprived they are but by which items they prioritise protecting, with adult-oriented items such as holidays and home repairs typically sacrificed before child-specific items such as warm coats. If households in this paper's sample sit in different latent classes with different binding items, a common income shock relaxes a different constraint for each one, which is consistent with a real aggregate improvement carrying no single robust item-level signal.
 
\subsection{Policy implications}
\label{sec:disc-policy}
 
Three implications follow from these results, for SCP specifically and for child income transfers more broadly.
 
First, the diffuse item-level effect above is itself evidence for cash over in-kind transfers as the right instrument here. An earmarked benefit tied to one hardship, such as a school-meals or winter-fuel voucher, can only ever move the specific item it targets. SCP's unrestricted cash instead let each household direct the payment toward whichever constraint it faced, consistent with an aggregate improvement that shows up nowhere in particular at the item level. A policy design that expected, or required, a detectable effect on one specific measured hardship would have misread this pattern as evidence the payment was not working.
 
Second, the gap between this paper's OLS and doubly robust estimates (Section~\ref{sec:disc-causal}) is a methodological lesson for future evaluations. Wales's Cynnal pilot (Section~\ref{sec:introduction}) will eventually face the same evaluation choice this paper faced for SCP: a simple OLS specification would likely overstate its effect relative to a doubly robust alternative. Evaluations of Cynnal, and of any future UK child payment, should report both rather than treating an OLS estimate as sufficient on its own.
 
Third, these results speak to policy beyond Scotland. The UK Government's removal of the two-child Universal Credit limit and its forthcoming Child Poverty Strategy (Section~\ref{sec:introduction}) both extend unrestricted cash-based support to a UK-wide population that includes England's children, the comparison group this paper uses as a counterfactual. This paper's evidence that a modest, unconditional cash payment can move material deprivation, without a traceable single-item mechanism, is directly relevant to how that UK-wide policy should be evaluated once it takes effect.